% Source: https://artofproblemsolving.com/community/c3h2037701_logical_problem
% Stable identity: aops-topic-2037701-q01
Paint used is proportional to painted area. If the large cube has edge \(L\) and the small cubes have edges \(s_1,\dots,s_k\), the fraction of paint on the outside is
\[
\dfrac{6L^2}{6(s_1^2+\cdots+s_k^2)}=\dfrac{L^2}{s_1^2+\cdots+s_k^2}.
\]

\emph{Equal cubes do not work.} If the large cube is made of \(n^3\) equal cubes of edge \(\dfrac {L}{n}\), the fraction is \(\dfrac{L^2}{n^3\cdot \dfrac{L^2}{n^2}}=\dfrac{1}{n}\), and \(\dfrac{1}{n}=\dfrac{12}{100}\) gives \(n=\dfrac{25}{3}\), not an integer.

\emph{Unequal cubes do work.} Take a \(3\times3\times3\) cube made of \(27\) unit cubes, and cut \(24\) of those unit cubes into \(27\) little cubes of edge \(\dfrac{1}{3}\) each.

\begin{asy}
usepackage("amsmath");
size(5cm);
defaultpen(fontsize(10pt));
real s=1;
// big 3x3x3 cube in oblique projection, one unit cube subdivided into thirds
pair ex=(1,0), ey=(0.45,0.35), ez=(0,1);
pair P(real x,real y,real z){return x*ex+y*ey+z*ez;}
for(int i=0;i<=3;++i){draw(P(i,0,0)--P(i,0,3),gray); draw(P(0,0,i)--P(3,0,i),gray);
 draw(P(3,i,0)--P(3,i,3),gray); draw(P(3,0,i)--P(3,3,i),gray);
 draw(P(0,i,3)--P(3,i,3),gray); draw(P(i,0,3)--P(i,3,3),gray);}
draw(P(0,0,0)--P(3,0,0)--P(3,3,0)--P(3,3,3)--P(0,3,3)--P(0,0,3)--cycle,linewidth(1));
draw(P(0,0,3)--P(3,0,3)--P(3,0,0));draw(P(3,0,3)--P(3,3,3));
for(int i=1;i<3;++i){draw(P(i/3,0,2)--P(i/3,0,3),blue+linewidth(0.4)); draw(P(0,0,2+i/3)--P(1,0,2+i/3),blue+linewidth(0.4));
 draw(P(i/3,0,3)--P(i/3,1,3),blue+linewidth(0.4)); draw(P(0,i/3,3)--P(1,i/3,3),blue+linewidth(0.4));}
label("$3$",P(1.5,0,0),S);
label("\mbox{}",(min(currentpicture,user=true)+max(currentpicture,user=true))/2,black); // keeps the text after the figure black when LaTeX typesets the labels inline
\end{asy}

Now the large cube consists of \(3\) unit cubes and \(24\cdot27=648\) cubes of edge \(\dfrac{1}{3}\). The total painted area is
\[
6\left(3\cdot1^2+648\cdot\dfrac{1}{9}\right)=6(3+72)=450,
\]
while the surface of the large cube has area \(6\cdot 3^2=54\). Since \(\dfrac{54}{450}=\dfrac{12}{100}\), exactly \(12\%\) of the paint is on the surface.

\[
\boxed{\text{Yes (with cubes of different sizes; impossible if all the small cubes are equal).}}
\]
